\section{Result package and reproduction limits}
\label{app:release}

Rebuilding this note's PDFs assembles documents; it does not reproduce
scientific results. The supporting internal records, including the digests
of the fixed identities, are kept in Ref.~\cite{NoteRecords2026}. This appendix distinguishes the existing local replay package from
the broader event-level and covariance release, states which identities are fixed,
and states the remaining reproduction limits.

% DECISION-20261021-joseph-da-and-note-records item 2, director ruling Q2 and review F6: the deposit snapshot hold;
% only the subject changed, and the clause from "refer to" on is byte-identical to app_provenance.tex lines 9-10
The code paths and file names in this appendix and in App.~\ref{app:codebase} refer to the snapshot of the code repository, \repolink,
identified at deposit: \depositsnapshot.

\paragraph{Release status.} A local release candidate, RC7, exists with a
README, checksums and a replay verifier~\cite{ReleaseCandidate2026}.
% Added 2026-10-21 with the shortened article Data availability (DECISION-20261021-joseph-da-and-note-records item 1;
% DA-proposal row 17, review D-2): RC7-README.md:75-78, :105-106 (fresh environment, empty directory, verify_rc.py
% VERIFY: PASS). RC7's own verifier, not an independent replay; the independent empty-directory replay was of RC1 (below).
RC7's own verifier, run on the packaged files in an empty directory,
reproduces every value recorded in the package and every figure number that
it checks. Its
approximately \SI{16}{MB} sufficient-input package reproduces the joint-test
$p$ values, decisions, robustness and power results, the lost-seed resolution
under reading~(a),
coarse-projection comparison and fine-grid $L_2$ ratios. It also regenerates,
from packaged arrays, figures that the article cites from this note: the
$(\Eavail,W)$ excess map (Fig.~\ref{fig:eavailWexcess}), the null
distributions of the joint tests (Fig.~\ref{fig:joint5d-nulls}), and the
one-dimensional generator projections.
% plain-language pass 2 (DECISION-20261018-joseph-plain-language-pass2): note home for the paper's
% Data availability sentence (pass-1 paper P6), from main_paper.tex "The package was assembled before ...".
% Source (as the paper's comment): Fig. 1 is stageF/rebuild/outputs/fig1/fig1_validation.pdf (82665808...),
% drawn by rebuild/scripts/make_fig1_q.py with the release's make_figs.py (9c4dbf24...) from cov2d_Q_205
% (rebuild/outputs/outA/cov2d_Q_205.json.gz) and budget_Q.
RC7 was assembled before the two-dimensional central value was changed to
LightGBM, so its two-dimensional curve is the exact-split one. The article's
two-dimensional validation figure and its numbers were instead regenerated
separately, from the committed LightGBM values and covariance, with RC7's
plotting code extended to draw this analysis's band.
It has not been
deposited or released under a publication tag. The schema below remains a
specification for the broader central-value and covariance products, not a
claim that RC7 contains them all. Fixed identities do not supply missing
scientific qualification.

\subsection{What today's reader can and cannot recompute}
\label{app:release-scope}

Public ancillary files supply the published inputs for comparisons to
Ref.~\cite{MINERvA:2021owq}'s central values and covariance, standardized
differences, tension anatomy, and MINERvA~Tune~v1 model comparison.

They do not reproduce this analysis's unfolded 2D uncertainty construction.
RC7 supports replay of the higher-dimensional inference from sufficient
inputs, but not extraction from the archived pseudo-experiment products,
event-level unfolding, or the quoted closure, bias and covariance-status
measurements. Nor does it reproduce the 2D uncertainty and coverage figures,
the low-recoil comparisons, the recoil-response sensitivity, or the proxy
convergence rate of joint-test condition~(i). The multiplier scan of the joint
tests is recomputable from RC7's sufficient inputs with the scan code of
Ref.~\cite{JointTestKappa2026}, but it is not part of RC7's own verification.
The pseudo-experiment products are archived (about \SI{2.3}{GB}). An earlier
version of the package was also replayed from an empty directory by an
independent reader. Event processing requires the earlier production of MINERvA
\href{https://doi.org/10.15484/3022562}{open-data AnaTuples}, event-loop outputs, and the
generator samples used here.

The repository now records SHA256 identities for all 2,374 inventoried
AnaTuple input files~\cite{AnaTuplePreservation2026}; RC7 itself still carries
the earlier names-and-sizes inventory. The dated 2026-10-10 preservation
record verifies copies of all 1,885 data, flux and parameter files on
collaboration storage. The 489 simulation AnaTuples are identified by hash
but have no preservation copy in that record. Internal preservation and
complete identities are distinct from a public deposit and do not make the
full event-level analysis reproducible from the small candidate package.

\subsection{Release schema}
\label{app:release-schema}

The broader result package must contain every item below. A $\checkmark$ marks identities or
values already fixed (bin edges: \S\ref{app:release-edges}; digests: Ref.~\cite{NoteRecords2026}, identities already fixed); the rest are packaging.

\begin{enumerate}[itemsep=2pt]
  \item \textbf{Reporting bin edges and units, per axis}
        $\checkmark$: $\pT$ and $\pPar$ in \si{GeV/c}, $\Eavail$, $q_{3}$ and
        $W$ in \si{GeV}, with each wide terminal catch bin identified.
  \item \textbf{The flattened-cell ordering} $\checkmark$: the dense
        $14\times16\times7\times7\times6$ grid in C order, with the reported
        subset supplied as an explicit index array, never inferred from a
        positivity cut.
  \item \textbf{Central values}: the five-axis array and every projection
        quoted, each with its own digest, on the density scale in
        \si{cm^2/nucleon} divided by the product of its kept axes' units.
  \item \textbf{Covariances} $\checkmark$: the adopted five-axis object and its
        $(\Eavail,W)$ projection, with the algebra of
        Eqs.~\eqref{eq:zassembly}--\eqref{eq:zinflation}, the declared
        centering variant, and the band partition.
  \item \textbf{Masks} $\checkmark$: the support mask, the pinned mask (bins
        where the inflation denominator vanished and $g$ was declared $1$),
        and the unreported-cell mask.
  \item \textbf{Projection matrices} $\checkmark$: $M$ itself, with its
        content digest, shape, and explicit identification of its entries as
        dropped axes' bin-width products, needed to specify destination units.
  \item \textbf{Generator predictions}: the four truth-level $(\Eavail,W)$
        predictions (GENIE-CV, GENIE$+$Valencia~2p2h, NuWro, GiBUU) and the
        MINERvA~Tune~v1 ancillary model, as arrays on the release binning
        rather than figure curves, each with generator version,
        configuration, target and flux. GENIE and NuWro require the
        flux-repaired predictions of 2026-09-27 (per-event reweight to the
        analysis PPFX flux plus the 50--\SI{100}{GeV} supplements); GiBUU
        requires its reweight to the same flux convention (Ref.~\cite{NoteRecords2026},
        \texttt{flux-repair}; \S\ref{sec:3d-models}).
        % KNOWN_ISSUES.md rows 82--83
  \item \textbf{Provenance}: each object's producing revision, job, and
        digest of bytes read back from the closed file. Only the assembling
        revision and job of $C_{Z}$, job of $C_{EW}$ and read-back row-index
        digest are fixed (Ref.~\cite{NoteRecords2026}, identities already fixed). Other objects lack
        recorded producing revisions and jobs; the assembled package carries
        none.
  \item \textbf{One minimal worked example}: \S\ref{app:release-example}.
  \item \textbf{The four measurements} that travel with the adopted covariance
        (\S\ref{sec:eavailw-adopted}) must appear \emph{in} the package readers
        will quote, beyond their inclusion in this note.
\end{enumerate}

\subsection{Bin edges}
\label{app:release-edges}

\noindent The bin edges are the one thing a reader
cannot recover from a digest:
\begin{itemize}[nosep,leftmargin=2em]
  \item $\pT$, \si{GeV/c}:
        $0$, $0.07$, $0.15$, $0.25$, $0.33$, $0.40$, $0.47$, $0.55$, $0.70$,
        $0.85$, $1.00$, $1.25$, $1.50$, $2.50$, $4.50$ (14 bins).
  \item $\pPar$, \si{GeV/c}:
        $1.5$, $2$, $2.5$, $3$, $3.5$, $4$, $4.5$, $5$, $6$, $7$, $8$, $9$,
        $10$, $15$, $20$, $40$, $60$ (16 bins).
  \item $\Eavail$, \si{GeV}: $0$, $0.1$, $0.2$, $0.4$, $0.8$, $1.5$, $3.0$,
        $100$ (7 bins).
  \item $q_{3}$, \si{GeV}: $0$, $0.2$, $0.4$, $0.6$, $0.8$, $1.2$, $2.0$,
        $100$ (7 bins).
  \item $W$, \si{GeV}: $0$, $1.1$, $1.4$, $1.8$, $2.2$, $3.0$, $100$
        (6 bins).
\end{itemize}
The last bin on $\Eavail$, $q_{3}$ and $W$ is a catch bin to \SI{100}{GeV}, and
the $\pT\times\pPar$ grid carries the 19 unreported high-$\pT$, low-$\pPar$
acceptance cells of \S\ref{sec:method-binning}.

\subsection{One minimal projection example}
\label{app:release-example}

Readers should be able to run this complete computation to check the released
ordering, mask and weights.

\begin{enumerate}[itemsep=2pt]
  \item Load the five-axis central values $\mathbf{x}$ and the reported-cell
        index array. Assert $\mathbf{x}$ has \num{65856} entries and the index
        array \latBins.
  \item Load $M$ (shape $\zprojCells\times\latBins$) and $C_{Z}$ (shape
        $\latBins\times\latBins$).
        The five-axis covariance used in this note is the file \texttt{z-cv.npz}, with \texttt{sha256}
        \zcvDigest\ldots\ and $\sqrt{\Tr C}=\SI{5.6742e-38}{}$, so a loaded $C_{Z}$ can be checked against both.
  \item Form $\mathbf{y}=M\mathbf{x}_{\text{reported}}$ and reshape to
        $7\times6$ \textbf{in C order}. This must reproduce the
        $(\Eavail,W)$ central-value map of Fig.~\ref{fig:eavailWexcess}. The
        F-order control establishes the ordering: C order agrees with the
        figure to a relative $1.7\times10^{-16}$; F order disagrees at relative
        $1.0$.
  \item Form $C_{EW}=M\,C_{Z}\,M^{\top}$ and check
        $\sqrt{\Tr C_{EW}}=\SI{4.4552e-39}{}$ and exact symmetry
        ($\max|C-C^{\top}|=0$). Do not check a rank. The spectrum decays
        smoothly with no plateau. The cutoff sets the retained count: 26 at
        $10^{-6}$, 36 at the projector's hardcoded $10^{-12}$, 41 at
        \texttt{numpy.linalg.matrix\_rank}'s default, 42 with no cutoff.
        A ``rank should be 36'' check using NumPy defaults would return 41 and
        falsely flag corruption. Check that
        $|\lambda_{\min}|/\lambda_{\max}$ is of order $10^{-15}$, the edge of
        double precision. Treat the object as ill-conditioned, without a
        definite rank. Do not check the sign of $\lambda_{\min}$, which sits
        at the noise floor and carries no information. A second host
        reproduces $\lambda_{\min}$ to only five significant digits.
  \item Read the per-cell relative uncertainty as
        $\sqrt{(C_{EW})_{rr}}/y_{r}$. Because $M$ carries the dropped axes'
        widths, $\mathbf{y}$ and $C_{EW}$ are on the \emph{density} scale. Converting cell $r$ to an integrated cross section requires
        multiplying by that cell's own $\Delta\Eavail\,\Delta W$, and the
        catch bins make that factor large.
  \item Do \emph{not} invert $C_{EW}$ without a declared regularization: its
        spectrum runs continuously down to the double-precision noise floor, so
        the inverse is set by the regularizer rather than by the data, and
        \S\ref{sec:eavailw-adopted}'s first measurement is the reason no
        significance is computed from it at all.
\end{enumerate}

\paragraph{Undeclared acceptance field.} The
projection receipt records \texttt{acceptance\_question: UNDECLARED}: whether
an acceptance correction is owed at the destination binning was undecided
when the projection was built. The package must retain this field and value
until a decision is made, because a silent default would let the data format make a
scientific choice.
